\documentclass{article}

% Language setting
% Replace `english' with e.g. `spanish' to change the document language
\usepackage[english]{babel}

% Set page size and margins
% Replace `letterpaper' with `a4paper' for UK/EU standard size
\usepackage[a4paper,top=2cm,bottom=2cm,left=3cm,right=3cm,marginparwidth=1.75cm]{geometry}

% Useful packages
\usepackage{amsmath}
\usepackage{graphicx}
\usepackage[colorlinks=true, allcolors=blue]{hyperref}

\title{Statistics Course}
\author{Wisam Yassar Mahardika}
\date{April 21\textsuperscript{th} 2026}

\begin{document}
\maketitle
\clearpage
\tableofcontents
\clearpage
\section{Key Words}
\begin{enumerate}
	\item \textbf{Data}: Data is any observations have been collected.
	\item \textbf{Statistics}: Collecting, analyzing, summarizing, interpreting and drawing conclusions from data. \textbf{Note}: \textit{data is useless without doing statistics}
	\item \textbf{Population}: The complete set of \textit{elements} being studied.
	\item \textbf{Sample}: Some subset of a population. \textbf{Note}: \textit{statistics usually deals with samples}
	\item \textbf{Census}: Collecting from every member of a population.
	\item \textbf{Parameter}: A characteristic of a population. \textit{(PP / Population Parameter).}
\item \textbf{Statistic}: A characteristic of a sample. \textit{(SS / Statistic Sample).}
\end{enumerate}
$ \rightarrow $ If you take a sample it must be collected randomly otherwise it's going to have bias on it.
\subsection{Two Types of Data}
\begin{enumerate}
\item \textbf{Quantitative}: Numerical data. \textit{e.g} height, weight, wages, temperature. \textbf{Note}: mathematical operations are \textit{meaningful.}
	\item \textbf{Qualitative(Categorical)}: Non-numeric data. \textit{e.g} gender, race, religion. \textbf{Note}: mathematical operations are \textit{meaningless.}
\end{enumerate}
\subsubsection{Two Types of Quantitative Data}
\begin{enumerate}
	\item \textbf{Discrete Data}: Countable or finite. \textit{e.g} \# of eggs, dice. \textbf{Note}: usually a count.
	\item \textbf{Continous Data}: Infinite numbers of possible values \textit{(not countable)}. \textit{e.g} temperature. \textbf{Note}: usually a measurement.
\end{enumerate}
\subsection{4 Levels of Measurement}
\begin{enumerate}
	\item \textbf{Nominal}: Categories \textbf{NOT} ordered. \textit{e.g} religion.
	\item \textbf{Ordinal}: Can be ordered. \textbf{Note}: differences are meaningless. \textit{e.g} rank, colors(spectrum).
	\item \textbf{Interval}: Can be ordered, differences are meaningful, no \textit{``Natural Zero``.} \textit{e.g} temperature.
	\item \textbf{Ratio}: Just like interval but with a \textit{``Natural Zero``.} \textit{e.g} amount of money.
\end{enumerate}
\end{document}
